\documentclass[11pt]{article}
\usepackage[utf8]{inputenc}
\usepackage{amsmath, amssymb, amsthm}
\usepackage[margin=1in]{geometry}

\theoremstyle{definition}
\newtheorem{definition}{Definition}
\newtheorem{theorem}{Theorem}

\theoremstyle{remark}
\newtheorem{remark}{Remark}

\DeclareMathOperator{\tr}{tr}

\title{Lecture 3: Stochastic Integration in Hilbert Space}
\date{30.10.2025}
\author{Tien Anh Vu}
\begin{document}
\maketitle

%% =============================
%% §3 Stochastic integration in Hilbert space
%% =============================
\section*{\S 3 Stochastic Integration in Hilbert Space}

%% =============================
%% 1.3.1 Martingales in Banach space
%% =============================
\subsection*{1.3.1 Martingales in Banach space}

Let $(E, \|\cdot\|)$ be a (real) Banach space.

%% =============================
%% Bochner integral
%% =============================
\subsubsection*{Bochner integral (E-valued integration)}

% The Bochner integral extends the Lebesgue integral to functions
% taking values in a Banach space E.

Given a finite measure $\mu$ on a probability space $(\Omega, \mathcal{A})$,
we want to define
\[
    \int f\,d\mu \;\in\; E.
\]

\begin{remark}
The Bochner integral is the natural generalisation of the Lebesgue integral to functions taking values in a Banach space~$E$.  Just as the Lebesgue integral is built by first defining the integral for simple (step) functions and then extending by density, the Bochner integral follows the same strategy but for $E$-valued functions.  The key challenge in infinite dimensions is that one must distinguish between different notions of measurability (strong vs.\ weak), which coincide in finite-dimensional spaces but can differ in general Banach spaces.
\end{remark}

%% =============================
%% Step 1: definition for step functions
%% =============================
\paragraph{Step functions.}
% We first define the integral for simple (step) functions,
% then extend by taking limits.

For measurable \textbf{vector-valued} functions $f \colon \Omega \to E$,
given a step function of the form
\[
    f \;=\; \sum_{k=1}^{n} x_k\,\mathbf{1}_{A_k},
    \qquad x_k \in E,\quad A_k \in \mathcal{A},\quad
    A_k \cap A_\ell = \emptyset, \tag{$*$}
\]
define
\[
    \int f\,d\mu
    \;:=\;
    \sum_{i=1}^{n} \mu(A_k)\,x_k
    \qquad \text{(linear combination of vectors $x_k$)}.
\]

\begin{remark}
The integral of a step function is a finite linear combination of vectors $x_k$ weighted by the measures of the corresponding sets~$A_k$.  This is completely analogous to the Lebesgue integral of real-valued simple functions.  The requirement that the sets $A_k$ are pairwise disjoint ensures that the definition is unambiguous: each $\omega \in \Omega$ belongs to at most one $A_k$, so $f(\omega) = x_k$ is uniquely determined.
\end{remark}

%% =============================
%% Bochner inequality
%% =============================
\paragraph{Bochner inequality.}
% This is the key estimate that allows extending
% the integral to limits of step functions.

We want to take suitable limits. The \textbf{Bochner inequality} gives:
\[
    \left\|\int f\,d\mu\right\|
    \;=\;
    \left\|\sum_{k=1}^{n} \mu(A_k)\,x_k\right\|_E
    \;\le\;
    \sum_{k=1}^{n} \|x_k\|_E\;\mu(A_k)
    \;=\;
    \int_\Omega \|f\|_E\,d\mu.
\]

% For a step function with disjoint A_k, the norm function
% ||f|| is itself a simple real-valued function:
For disjoint $A_k$, the norm of $f$ is the real-valued step function:
\[
    \|f\|(\omega)
    \;=\;
    \sum_{k} \|x_k\|\;\mathbf{1}_{A_k}(\omega),
\]
and the right-hand side is just the ordinary \textbf{Lebesgue integral}
of $\|f\|_E$.

\begin{remark}
The Bochner inequality is the key estimate that enables extending the integral from step functions to general Bochner-integrable functions.  If a sequence of step functions $f_n \to f$ converges in the $L^1$ norm, then by the Bochner inequality the integrals $\int f_n\,d\mu$ form a Cauchy sequence in~$E$, and since $E$ is complete (being a Banach space), the limit exists and defines $\int f\,d\mu$.  This is exactly the same density argument used to extend the Lebesgue integral from simple functions to $L^1$.
\end{remark}

%% =============================
%% Extension to general Bochner-integrable functions
%% =============================
\paragraph{Extension to all integrals.}
% The Bochner inequality allows us to extend the integral
% from step functions to the closure in L^1 norm.

The Bochner inequality allows us to extend $\int f\,d\mu$ to all functions
\[
    f \colon \Omega \to E
    \quad \text{satisfying:}
\]
\begin{itemize}
    \item $f$ is \textbf{strongly measurable}, i.e.,
          \[
              f = \lim_{n \to \infty} f_n
              \qquad \mu\text{-a.e.\ (in $E$)},
          \]
          where $f_n$ are of type $(*)$,
          % i.e., f can be approximated pointwise a.e. by step functions

    \item $\displaystyle\int_\Omega \|f\|_E\,d\mu < \infty$.
          % The E-norm of f must be Lebesgue integrable
\end{itemize}

%% =============================
%% Remark: strong vs weak measurability
%% =============================
\begin{remark}
In general, it is \textbf{not} true that if
$f \colon \Omega \to E$ is $\mathcal{A}/\mathcal{B}(E)$-measurable,
this implies that $f$ can be approximated by step functions of type $(*)$.
% Borel measurability alone is insufficient in infinite-dimensional spaces.
Rather, such a function $f$ is only \textbf{weakly measurable}, i.e.,
for any $\ell \in E'$ (= all continuous linear functionals on $E$),
the real-valued map $\omega \mapsto \ell(f(\omega))$ is measurable.
% The distinction between strong and weak measurability is crucial
% in infinite-dimensional Banach spaces; they coincide in separable spaces
% by the Pettis measurability theorem.
\end{remark}

%% =============================
%% Pettis measurability theorem
%% =============================
\begin{theorem}[Pettis measurability theorem]
A necessary and sufficient condition for
$f \colon \Omega \to E$ to be \textbf{strongly measurable} is that $f$ is
\textbf{weakly measurable} and \textbf{separably valued}
(i.e., $f(\Omega) \subseteq E$ is separable).
\end{theorem}

\begin{remark}
In particular, if $E$ is separable
(e.g.\ $\mathbb{R}^d$, $L^p(\mathbb{R}^d)$, etc.),
then every weakly measurable mapping $f \colon \Omega \to E$ is
also \textbf{strongly measurable}.
% This is because in a separable Banach space, the Borel sigma-algebra
% coincides with the sigma-algebra generated by continuous linear functionals.
\end{remark}

%% =============================
%% Notation: Bochner-Lebesgue spaces
%% =============================
\paragraph{Notation.}
% Bochner--Lebesgue spaces are the E-valued analogues
% of the classical L^p spaces.
We write
\[
    L^p(\Omega, \mathcal{A}, \mu;\, E),
    \qquad
    L^p(\Omega, \mathcal{A}, \mu,\, E)
\]
for the space of all Bochner-measurable functions
$f \colon \Omega \to E$ with
$\int_\Omega \|f\|_E^p\,d\mu < \infty$.

\begin{remark}
The Bochner--Lebesgue spaces $L^p(\Omega, \mathcal{A}, \mu;\, E)$ are the natural framework for random variables and stochastic processes taking values in a Banach space~$E$.  For $p = 1$, the space consists of all Bochner-integrable functions.  For $p = 2$ and $E$ a Hilbert space, the space $L^2(\Omega, \mathcal{A}, \mu;\, H)$ is again a Hilbert space, which is essential for developing the $L^2$ theory of stochastic integration.
\end{remark}

%% =============================
%% Properties of Bochner's integral
%% =============================
\subsection*{Properties of Bochner's integral}

\begin{itemize}
    \item[(i)] \textbf{Bochner inequality:}
    \[
        \left\|\int f\,d\mu\right\|
        \;\le\;
        \int \|f\|\,d\mu.
    \]
    % This is the E-valued analogue of the triangle inequality for integrals.

    \item[(ii)] \textbf{Linearity:} If $L \colon E \to F$ is a
    \textbf{continuous, linear} map (where $F$ is also a Banach space), then
    \[
        L\!\left(\int f\,d\mu\right)
        \;=\;
        \int L(f)\,d\mu.
    \]
    % Continuous linear operators commute with the Bochner integral,
    % which is an interchange of L and the integral.

    \item[(iii)] If $f \in C\bigl([a,b];\, E\bigr)$,
    then
    \[
        f(t) - f(s) \;=\; \int_s^t f'(r)\,dr,
    \]
    where the right-hand side is a \textbf{Bochner integral}.
    % This is the fundamental theorem of calculus for E-valued functions.

\end{itemize}

\begin{remark}
Property~(iii) is the \textbf{fundamental theorem of calculus} for Banach-space-valued functions.  It asserts that if $f \colon [a,b] \to E$ is continuously differentiable, then $f$ can be recovered from its derivative by Bochner integration, exactly as in the scalar case.  This result is used extensively in the theory of $C_0$-semigroups and evolution equations, where solutions are often expressed as integrals of operator-valued functions.
\end{remark}

%% =============================
%% Conditional expectation in Banach spaces
%% =============================
\subsection*{Conditional expectation in Banach spaces}

% Here (\Omega, \mathcal{A}, P) is a probability space.

\begin{theorem}[Theorem 1.16 --- Conditional expectation]
Let $X$ be a random variable $X \colon \Omega \to E$
with $X \in L^1(\Omega, \mathcal{A}, P;\, E)$,
i.e., $X$ is a Bochner-integrable $E$-valued random variable.
Let $\mathcal{A}_0 \subset \mathcal{A}$ be a sub-$\sigma$-algebra of
$\mathcal{A}$.  Then there exists a $P$-a.s.\ uniquely determined
\[
    X_0 \;\in\; L^1(\Omega, \mathcal{A}, P;\, E)
\]
such that
\[
    \int_A X_0\,dP
    \;=\;
    \int_A X\,dP
    \qquad
    \forall\, A \in \mathcal{A}_0.
\]
$X_0$ is called the \textbf{conditional expectation} of $X$ given
$\mathcal{A}_0$, and we write
$X_0 \equiv E(X \mid \mathcal{A}_0)$.
\end{theorem}

\begin{remark}
Bochner's inequality implies the following estimate for
conditional expectations:
\[
    \bigl\|E(X \mid \mathcal{A}_0)\bigr\|_E
    \;\le\;
    E\!\bigl(\|X\|_E \;\big|\; \mathcal{A}_0\bigr).
\]
% The norm of the conditional expectation is controlled by the
% conditional expectation of the norm.
% In the Hilbert space setting, it is not possible to require
% monotonicity of conditional expectations (as in the real-valued case),
% because there is no natural ordering on a Hilbert space.
In a Hilbert space, it is \textbf{not} possible to require
monotonicity ($\ge$) for conditional expectations,
since there is no natural partial order on $H$.
\end{remark}

%% =============================
%% Proof sketch of Theorem 1.16
%% =============================
\paragraph{Proof (sketch).}
Start with a step function
$X = \sum_{k=1}^{n} x_k\,\mathbf{1}_{A_k}$
(pairwise disjoint $A_k$).  Then
\[
    E(X \mid \mathcal{A}_0)
    \;=\;
    \sum_{k=1}^{n} x_k\;E\!\bigl(\mathbf{1}_{A_k}
    \mid \mathcal{A}_0\bigr).
\]
% Each term uses the real-valued conditional expectation
% of an indicator function, which is well-defined.

Hence,
\[
    \bigl\|E(X \mid \mathcal{A}_0)\bigr\|
    \;\le\;
    \sum_{k=1}^{n} \|x_k\|\;
    E\!\bigl(\mathbf{1}_{A_k} \mid \mathcal{A}_0\bigr)
    \;=\;
    E\!\bigl(\|X\| \mid \mathcal{A}_0\bigr).
\]
% The general case follows by approximating an arbitrary
% Bochner-integrable function by step functions and passing
% to the limit using converging sequences.
Generalize by using converging sequences.

\begin{remark}
The extension from step functions to general $E$-valued random variables follows the same approximation scheme as in the scalar case: approximate $X$ by step functions $X_n$, apply the conditional expectation to each~$X_n$ (which is explicit), and pass to the limit.  The Bochner inequality ensures that the conditional expectation operator is a \textbf{contraction} on $L^1(\Omega, \mathcal{A}, P;\, E)$: if $X_n \to X$ in $L^1$, then $E(X_n \mid \mathcal{A}_0) \to E(X \mid \mathcal{A}_0)$ in $L^1$, guaranteeing existence and uniqueness of the limit.
\end{remark}

%% =============================
%% Definition: E-valued martingale
%% =============================
\subsection*{Martingales in Banach spaces}

\begin{definition}[$E$-valued martingale]
An $E$-valued stochastic process $(M_t)_{t \ge 0}$ is called an
$(\mathcal{A}_t)$-\textbf{martingale}
(with respect to the filtration $(\mathcal{A}_t)$) if:
\begin{itemize}
    \item[(i)] $E\!\bigl(\|M_t\|\bigr) < \infty$
          \quad (Bochner integrability),
    \item[(ii)] $M_t$ is $\mathcal{A}_t$-adapted,
    \item[(iii)] $E(M_t \mid \mathcal{A}_s) = M_s$,
          \quad $0 \le s \le t$.
\end{itemize}
\end{definition}

\begin{remark}
The definition directly generalises scalar-valued martingales to the Banach space setting.
Condition~(i) ensures that the Bochner integral (and hence the conditional expectation in~(iii)) is well-defined; it replaces the classical requirement $E(|M_t|) < \infty$.
Condition~(ii) is the usual adaptedness: $M_t$ may only depend on the information available at time~$t$.
Condition~(iii) is the \textbf{martingale identity}, now interpreted via \textbf{Bochner conditional expectation} (Theorem~1.16).
The key subtlety in infinite dimensions is that the conditional expectation must be constructed using the Bochner integral theory, and one cannot rely on monotonicity arguments as in the real-valued case.
\end{remark}

\begin{remark}
If $(M_t)$ is an $E$-valued $(\mathcal{A}_t)$-martingale,
then for every $\ell \in E'$ (continuous linear functional),
$\bigl(\ell(M_t)\bigr)_{t \ge 0}$ is a \textbf{real-valued}
$(\mathcal{A}_t)$-martingale.
% This follows from property (ii) of the Bochner integral:
% continuous linear maps commute with the integral,
% hence with conditional expectations.
\end{remark}

%% =============================
%% Example: Q-Wiener process is a martingale
%% =============================
\paragraph{Example: $Q$-Wiener process.}
Let $(W_t)$ be a $Q$-Wiener process taking values in $U$
with filtration $(\mathcal{A}_t)$.
Then $(W_t)$ is an $(\mathcal{A}_t)$-martingale, since
$E\!\bigl(\|W_t\|_U^2\bigr) < \infty$.

% We verify the finite second moment using the canonical
% representation of the Q-Wiener process.

Indeed, using the \textbf{canonical representation} of
the Wiener process,
\begin{align*}
    E\!\left(\|W_t\|_U^2\right)
    &\;=\;
    E\!\left(
        \left\|\sum_{k=1}^{\infty}
        \sqrt{\lambda_k}\,\beta_k(t)\,e_k
        \right\|_U^2
    \right) \\[6pt]
    %% Apply Parseval's identity: ||x||^2 = sum |<x, e_k>|^2
    &\;=\;
    E\!\left(
        \sum_{k=1}^{\infty}
        \bigl\langle W_t,\, e_k \bigr\rangle_U^2
    \right)
    \qquad\text{(Parseval's identity)} \\[6pt]
    %% Each coordinate: <W_t, e_k> = sqrt(lambda_k) beta_k(t)
    &\;=\;
    E\!\left(
        \sum_{k=1}^{\infty}
        \lambda_k\,\beta_k^2(t)
    \right) \\[6pt]
    %% Interchange sum and expectation by Fubini/monotone convergence
    &\;\stackrel{\text{Fubini}}{=}\;
    \sum_{k=1}^{\infty} \lambda_k\;
    \underbrace{E\!\bigl(\beta_k^2(t)\bigr)}_{=\;t} \\[6pt]
    &\;=\;
    t\,\sum_{k=1}^{\infty} \lambda_k
    \;=\;
    t\;\tr(Q).
\end{align*}

\begin{remark}
Here $(e_k)$ is a CONS of $U$ consisting of eigenvectors of $Q$
with eigenvalues $\lambda_k$, and $(\beta_k(t))$ are independent
standard Brownian motions.
The interchange of summation and expectation is justified by
Fubini's theorem (or monotone convergence), since all terms
are non-negative.  The identity
$E(\beta_k^2(t)) = t$ is the variance of a standard Brownian motion.
\end{remark}

%% =============================
%% Martingale property of Q-Wiener process (verification)
%% =============================
\paragraph{Verification of the martingale property.}
% We check E(W_t | A_s) = W_s using independent increments.

For $s \le t$:
\[
    E(W_t \mid \mathcal{A}_s)
    \;=\;
    E\!\bigl(\underbrace{W_t - W_s}_{\text{independent of } \mathcal{A}_s}
    + W_s \;\big|\; \mathcal{A}_s\bigr)
    \;=\;
    \underbrace{E(W_t - W_s)}_{=\;0}
    + W_s
    \;=\;
    W_s.
\]
% The increment W_t - W_s is independent of A_s and has mean 0,
% so the Q-Wiener process is indeed a martingale.

Hence the $Q$-Wiener process is a martingale.

\begin{remark}
The verification uses two key properties: (1)~the increment $W_t - W_s$ is \textbf{independent} of $\mathcal{A}_s$ (by definition of the $Q$-Wiener process), and (2)~it has \textbf{mean zero} since $W_t - W_s \sim \mathcal{N}(0, (t-s)Q)$.  Independence allows us to factor the conditional expectation, and the zero-mean property gives the result.  This is the same argument as for scalar Brownian motion, but carried out in the Hilbert space~$U$ using the Bochner-integral formulation of conditional expectation.
\end{remark}

%% =============================
%% Doob's maximal inequality in Banach spaces
%% =============================
\subsection*{Doob's maximal inequality}

\begin{theorem}[Doob's maximal inequality]
Let $(M_t)$ be a right-continuous $E$-valued
$(\mathcal{A}_t)$-martingale.  Then for $p > 1$:
\[
    \left(
        E\!\left(\sup_{0 \le t \le T} \|M_t\|^p\right)
    \right)^{1/p}
    \;\le\;
    \frac{p}{p-1}\;
    \bigl(E\!\bigl(\|M_T\|^p\bigr)\bigr)^{1/p}
    \qquad \forall\, p > 1.
\]
\end{theorem}

\begin{remark}
Doob's maximal inequality controls the \textbf{supremum process}
$\sup_{0 \le t \le T} \|M_t\|$ in terms of the \textbf{terminal value}
$\|M_T\|$.  The constant $\frac{p}{p-1}$ diverges as $p \to 1^+$,
reflecting the fact that the inequality fails for $p = 1$
(one needs the weaker $L^1$ maximal inequality in that case).
For martingales, the key idea is that $\|M_t\|^p$ is a
\textbf{submartingale} (by Jensen's inequality applied to the
convex function $\|\cdot\|^p$), and the inequality then follows
from the classical Doob submartingale inequality.
This result is fundamental for proving convergence of stochastic
integrals and for establishing pathwise regularity of solutions to SPDEs.
\end{remark}

%% =============================
%% Space of L^p-bounded martingales
%% =============================
\paragraph{Space of $L^p$-bounded martingales.}
% Define the martingale M^p space with the supremum norm.

Let $\mathcal{M}^p := \bigl\{(M_t)_{t \in [0,T]} :
\text{$E$-valued continuous $(\mathcal{A}_t)$-martingale,
$L^p$-bounded}\bigr\}$ with norm
\[
    \|M\|_{\mathcal{M}^p}
    \;:=\;
    \sup_{0 \le t \le T}
    \bigl(E\!\bigl(\|M_t\|_E^p\bigr)\bigr)^{1/p}
    \;<\; \infty.
\]
This is a Banach space for $p > 1$ (Hilbert space if $p = 2$
and $E$ is a Hilbert space).

\begin{remark}
By Doob's maximal inequality, an equivalent norm on $\mathcal{M}^p$ is
\[
    \left(E\!\left(\sup_{0 \le t \le T}
    \|M_t\|_E^p\right)\right)^{1/p}.
\]
\end{remark}

%% =============================
%% 1.3.2 Stochastic integral
%% =============================
\subsection*{1.3.2 Stochastic integral}

Fix a $U$-valued $Q$-Wiener process $(W_t)$ with filtration
$(\mathcal{A}_t)$.  We want to define
\[
    \int \Phi(s)\,dW_s \;\in\; H.
\]

\begin{remark}
The construction of the stochastic integral $\int_0^T \Phi_s\,dW_s$ follows the same three-step programme as the classical It\^o integral in finite dimensions:
(1)~define the integral for \emph{elementary} (step) processes by a finite sum;
(2)~establish an \emph{isometry} (the Wiener--It\^o isometry) that controls the $L^2$ norm of the integral in terms of the integrand;
(3)~\emph{extend} the integral to all square-integrable predictable integrands by a density/completion argument.
The main new feature in infinite dimensions is that the integrand $\Phi_s$ is an \emph{operator} from $U$ to $H$ (not a scalar or a vector), and the isometry involves the \textbf{Hilbert--Schmidt norm} of $\Phi_s \circ \sqrt{Q}$ rather than the squared absolute value $|\phi_s|^2$ familiar from the scalar case.
\end{remark}

%% =============================
%% Step 1: Integration of elementary processes
%% =============================
\paragraph{Step 1: Integration of elementary processes.}

% An elementary process is a step function in time
% with values that are bounded linear operators U -> H.

An \textbf{elementary process} has the form:
\[
    \Phi(t)
    \;=\;
    \sum_{m=0}^{k-1} \Phi_m\;\mathbf{1}_{]t_m,\, t_{m+1}]}(t),
    \qquad
    0 = t_0 < t_1 < \cdots < t_k = T, \tag{$\star$}
\]
where each $\Phi_m$ is $\mathcal{A}_{t_m}$-measurable and
$\Phi_m \colon U \to H$ is a bounded linear operator,
i.e., $\Phi_m \in L(U, H)$.

% The stochastic integral of an elementary process is defined
% as a finite sum of martingale increments applied through the operators.

Then define
\[
    \int_0^T \Phi_s\,dW_s
    \;:=\;
    \sum_{m=0}^{k-1}
    \Phi_m\!\bigl(W_{t_{m+1} \wedge t} - W_{t_m \wedge t}\bigr)
    \;\in\; H.
\]
% Each term: the operator \Phi_m (from U to H) is applied to
% the Wiener increment W_{t_{m+1}} - W_{t_m} (which lives in U),
% yielding an element of H.

This yields a \textbf{linear mapping}:
\[
    I \colon \mathcal{E}
    \;\longrightarrow\;
    \mathcal{M}^2,
\]
where $\mathcal{E}$ denotes the space of all elementary processes
of type $(\star)$.

\begin{remark}
Each term $\Phi_m(W_{t_{m+1} \wedge t} - W_{t_m \wedge t})$ in the definition applies the \textbf{operator} $\Phi_m \colon U \to H$ to the Wiener increment $W_{t_{m+1} \wedge t} - W_{t_m \wedge t} \in U$, producing an element of $H$.  The adaptedness requirement ($\Phi_m$ is $\mathcal{A}_{t_m}$-measurable) ensures that the integrand is determined by information available \emph{before} the increment occurs --- this is the \textbf{non-anticipativity} condition that distinguishes It\^o integration from other stochastic calculi (e.g.\ Stratonovich).
\end{remark}

%% =============================
%% Step 2: Wiener--Itô isometry
%% =============================
\paragraph{Step 2: Wiener--Itô isometry.}

% If f takes values in a general Banach space, there is no
% Wiener--Itô isometry. This is specific to the Hilbert space setting.

\begin{remark}
If $\Phi$ takes values in a general Banach space,
there is \textbf{no} Wiener--Itô isometry.
The isometry requires $Q$ to be trace class and $\sqrt{Q}$
to be Hilbert--Schmidt.
\end{remark}

%% =============================
%% Hilbert--Schmidt operators
%% =============================
\paragraph{Hilbert--Schmidt operators.}

\begin{definition}[Hilbert--Schmidt operator]
An operator $L \in L(U, H)$ is called \textbf{Hilbert--Schmidt} if
\[
    \|L\|_{L_2(U,H)}^2
    \;:=\;
    \sum_{k=1}^{\infty} \|L\,e_k\|_H^2
    \;<\; \infty
\]
for all CONS $(e_k)$ of $U$.
\end{definition}

\begin{remark}
The definition is independent of the choice of CONS $(e_k)$.
% This can be verified by expanding in two different CONS
% and using Parseval's identity.
\end{remark}

\begin{remark}
Hilbert--Schmidt operators form a two-sided ideal in the algebra of bounded operators: if $A \in L_2(U, H)$ and $B \in L(H, V)$, $C \in L(V', U)$ are bounded, then $B \circ A$ and $A \circ C$ are also Hilbert--Schmidt.  The Hilbert--Schmidt norm satisfies $\|L\|_{\mathrm{op}} \le \|L\|_{L_2(U,H)}$, so every Hilbert--Schmidt operator is bounded, but the converse is false in infinite dimensions.  In finite dimensions ($U = \mathbb{R}^n$, $H = \mathbb{R}^m$), every linear operator is Hilbert--Schmidt and $\|L\|_{L_2}^2 = \sum_{i,j} L_{ij}^2$ equals the sum of squared matrix entries (the \textbf{Frobenius norm}).
\end{remark}

Let $L_2(U, H)$ be the (linear) space of all Hilbert--Schmidt operators
from $U$ to $H$.  Then $L_2(U, H)$ is again a (real) \textbf{Hilbert space}
with the scalar product
\[
    \langle A, B \rangle_{L_2(U,H)}
    \;:=\;
    \sum_{k=1}^{\infty}
    \langle A\,e_k,\; B\,e_k \rangle_H.
\]

%% =============================
%% Wiener--Itô isometry statement
%% =============================
\paragraph{Wiener--Itô isometry.}

Then the following \textbf{Wiener--Itô isometry} holds:
\[
    E\!\left(
        \left\|\int_0^t \Phi_s\,dW_s\right\|_H^2
    \right)
    \;=\;
    E\!\left(
        \int_0^t
        \bigl\|\Phi_s \circ \sqrt{Q}\,\bigr\|_{L_2(U,H)}^2
        \,ds
    \right).
\]

% In the M^2 norm, this reads:
In particular,
\[
    \left\|\int \Phi_s\,dW_s\right\|_{\mathcal{M}^{2}_{T}}^2
    \;=\;
    E\!\left(
        \int_0^T
        \bigl\|\Phi_s \circ \sqrt{Q}\,\bigr\|_{L_2(U,H)}^2
        \,ds
    \right),
\]
where the left-hand side is the norm in the space of continuous
$H$-valued, $\mathcal{A}_t$-adapted martingales.

\begin{remark}
The Wiener--It\^o isometry is the infinite-dimensional analogue of the classical It\^o isometry $E\!\bigl[\bigl(\int_0^T \phi_s\,dB_s\bigr)^2\bigr] = E\!\bigl[\int_0^T \phi_s^2\,ds\bigr]$ for scalar Brownian motion.  The key difference is that the integrand $\Phi_s$ is now an \emph{operator} from $U$ to $H$, and the squared integrand $|\phi_s|^2$ is replaced by the \textbf{Hilbert--Schmidt norm} $\|\Phi_s \circ \sqrt{Q}\|_{L_2(U,H)}^2$.  The appearance of $\sqrt{Q}$ reflects the spatial correlation of the driving noise: the composition $\Phi_s \circ \sqrt{Q}$ ``pre-whitens'' the noise, and the Hilbert--Schmidt norm measures the total energy injected into~$H$ per unit time.  This isometry is the foundation for extending the stochastic integral from elementary processes to all square-integrable predictable integrands via a density argument, exactly as the scalar It\^o integral is extended.
\end{remark}

%% =============================
%% Proof sketch of the Wiener--Itô isometry
%% =============================
\paragraph{Proof sketch.}
% Expand ||sum Phi_m Delta_m||^2 using bilinearity of the inner product.

We compute, writing $\Delta_m := W_{t_{m+1}} - W_{t_m}$:
\begin{align*}
    E\!\left(
        \left\|\int_0^T \Phi_s\,dW_s\right\|_H^2
    \right)
    &\;=\;
    E\!\left(
        \left\|\sum_{m=0}^{k-1}
        \Phi_m\,\Delta_m\right\|_H^2
    \right) \\[6pt]
    &\;=\;
    \sum_{m,n}
    E\!\bigl(
        \langle \Phi_m\,\Delta_m,\;
                \Phi_n\,\Delta_n \rangle_H
    \bigr).
\end{align*}

% For m ≠ n, the cross terms vanish due to the martingale property
% and independence of increments.
\begin{remark}
For $m \neq n$, the cross terms vanish:
$E\!\bigl(\langle \Phi_m\,\Delta_m,\;\Phi_n\,\Delta_n \rangle_H\bigr) = 0$.
This follows from the independence of the Wiener increments and
the adaptedness of $\Phi_m$.
\end{remark}

Hence only the diagonal terms $m = n$ survive.
Using the \textbf{canonical representation} of the Wiener process,
$\Delta_m = \sum_{k} \sqrt{\lambda_k}\;\bigl(\beta_k(t_{m+1}) -
\beta_k(t_m)\bigr)\,e_k$, we get:

\begin{align*}
    E\!\bigl(\langle \Phi_m\,\Delta_m,\;
             \Phi_n\,\Delta_n \rangle_H\bigr)
    &\;=\;
    \sum_{k,\ell}
    \sqrt{\lambda_k}\,\sqrt{\lambda_\ell}\;
    E\!\Bigl[
        \bigl\langle \Phi_m\!\bigl(e_k\bigr),\;
              \Phi_n\!\bigl(e_\ell\bigr)
        \bigr\rangle_H \\
    &\qquad\qquad\qquad \cdot
        \bigl(\beta_k(t_{m+1}) - \beta_k(t_m)\bigr)
        \bigl(\beta_\ell(t_{n+1}) - \beta_\ell(t_n)\bigr)
    \Bigr].
\end{align*}

\begin{remark}
When $m = n$, the expectation of the product of Brownian increments gives
$E\!\bigl((\beta_k(t_{m+1}) - \beta_k(t_m))(\beta_\ell(t_{m+1}) -
\beta_\ell(t_m))\bigr) = \delta_{k\ell}\,(t_{m+1} - t_m)$,
where $\delta_{k\ell}$ is the Kronecker delta. This is because the
$\beta_k$ are independent standard Brownian motions.
The sum then reduces to the Hilbert--Schmidt norm
$\|\Phi_m \circ \sqrt{Q}\|_{L_2(U,H)}^2$.
\end{remark}

%% =============================
%% Detailed computation of the cross and diagonal terms
%% =============================
\paragraph{Detailed computation (continued).}

% For m < n, use the tower property of conditional expectation and
% condition on the past sigma-algebra \mathcal{A}_{t_n}.
% The future increment \beta_\ell(t_{n+1})-\beta_\ell(t_n) has
% conditional mean zero given \mathcal{A}_{t_n}.

\noindent\textbf{Case $m < n$:} Starting from the cross term,
\begin{align*}
    &E\!\bigl(\langle \Phi_m\,\Delta_m,\;\Phi_n\,\Delta_n\rangle_H\bigr) \\[4pt]
    &\;=\;
    \sum_{k,\ell}\sqrt{\lambda_k}\,\sqrt{\lambda_\ell}\;
    E\!\Bigl(
        \bigl\langle \Phi_m(e_k),\;\Phi_n(e_\ell)\bigr\rangle_H\,
        \bigl(\beta_k(t_{m+1})-\beta_k(t_m)\bigr)\,
        \bigl(\beta_\ell(t_{n+1})-\beta_\ell(t_n)\bigr)
    \Bigr).
\end{align*}
Now apply the tower property with respect to $\mathcal{A}_{t_n}$:
\begin{align*}
    &E\!\Bigl(
        \bigl\langle \Phi_m(e_k),\;\Phi_n(e_\ell)\bigr\rangle_H\,
        \bigl(\beta_k(t_{m+1})-\beta_k(t_m)\bigr)\,
        \bigl(\beta_\ell(t_{n+1})-\beta_\ell(t_n)\bigr)
    \Bigr) \\[4pt]
    &\;=\;
    E\!\Bigl(
        E\!\Bigl[
            \bigl\langle \Phi_m(e_k),\;\Phi_n(e_\ell)\bigr\rangle_H\,
            \bigl(\beta_k(t_{m+1})-\beta_k(t_m)\bigr)\,
            \bigl(\beta_\ell(t_{n+1})-\beta_\ell(t_n)\bigr)
            \,\Big|\, \mathcal{A}_{t_n}
        \Bigr]
    \Bigr) \\[4pt]
    &\;=\;
    E\!\Bigl(
        \bigl\langle \Phi_m(e_k),\;\Phi_n(e_\ell)\bigr\rangle_H\,
        \bigl(\beta_k(t_{m+1})-\beta_k(t_m)\bigr)\,
        E\!\bigl[
            \beta_\ell(t_{n+1})-\beta_\ell(t_n)
            \,\big|\, \mathcal{A}_{t_n}
        \bigr]
    \Bigr) \\[4pt]
    &\;=\; 0.
\end{align*}
Here $\Phi_m$ and $\Phi_n$ are $\mathcal{A}_{t_n}$-measurable (since $m<n$ and the integrand is adapted), and $\beta_k(t_{m+1})-\beta_k(t_m)$ is also $\mathcal{A}_{t_n}$-measurable because $t_{m+1}\le t_n$. The future increment $\beta_\ell(t_{n+1})-\beta_\ell(t_n)$ is independent of $\mathcal{A}_{t_n}$ and has mean zero, so
\[
E\!\bigl[\beta_\ell(t_{n+1})-\beta_\ell(t_n)\mid \mathcal{A}_{t_n}\bigr]=0.
\]
Thus every summand is zero, and therefore the whole sum is zero. Similarly for $m > n$ by symmetry.

\begin{remark}
The vanishing of cross terms ($m \neq n$) is the fundamental \textbf{orthogonality property} of the stochastic integral and is the reason the It\^o integral is a martingale.
The argument uses the \textbf{tower property} of conditional expectation: for $m < n$, condition on $\mathcal{A}_{t_n}$ to separate the past increment $\beta_k(t_{m+1}) - \beta_k(t_m)$ (which is $\mathcal{A}_{t_n}$-measurable together with $\Phi_m$ and $\Phi_n$) from the future increment $\beta_\ell(t_{n+1}) - \beta_\ell(t_n)$ (which is independent of $\mathcal{A}_{t_n}$ with mean zero).
Factoring out the $\mathcal{A}_{t_n}$-measurable parts and using $E[\beta_\ell(t_{n+1}) - \beta_\ell(t_n) \mid \mathcal{A}_{t_n}] = 0$ yields zero.
This is the same mechanism that makes the classical scalar It\^o integral a martingale.
\end{remark}

% Now handle the case m = n in detail.
\noindent\textbf{Case $m = n$:}
\begin{align*}
    &\sum_{k,\ell}
    \sqrt{\lambda_k}\,\sqrt{\lambda_\ell}\;
    E\!\Bigl(
        \bigl\langle \Phi_m(e_k),\;\Phi_m(e_\ell)
        \bigr\rangle_H
        \cdot
        \bigl(\beta_k(t_{m+1}) - \beta_k(t_m)\bigr)
        \bigl(\beta_\ell(t_{m+1}) - \beta_\ell(t_m)\bigr)
    \Bigr) \\[6pt]
    %% Use independence: E(beta_k * beta_ell) = delta_{kl} * (t_{m+1} - t_m)
    &\;=\;
    \sum_{k}
    \lambda_k\;
    E\!\bigl(\|\Phi_m\,e_k\|_H^2\bigr)
    \cdot (t_{m+1} - t_m) \\[6pt]
    %% Recognize sqrt(lambda_k) e_k = sqrt(Q) e_k
    &\;=\;
    \sum_{k}
    E\!\bigl(
        \bigl\|\Phi_m\!\bigl(\sqrt{\lambda_k}\,e_k\bigr)
        \bigr\|_H^2
    \bigr)
    \cdot (t_{m+1} - t_m) \\[6pt]
    %% Since sqrt(Q) e_k = sqrt(lambda_k) e_k, this is ||Phi_m(sqrt(Q) e_k)||^2
    &\;=\;
    E\!\left(
        \sum_{k}
        \bigl\|\Phi_m \circ \sqrt{Q}\;e_k\bigr\|_H^2
    \right)
    \cdot (t_{m+1} - t_m) \\[6pt]
    %% The sum over k is the HS norm squared
    &\;=\;
    E\!\Bigl(
        \bigl\|\Phi_m \circ \sqrt{Q}\,\bigr\|_{L_2(U,H)}^2
    \Bigr)
    \cdot (t_{m+1} - t_m).
\end{align*}

% In the last step we used the definition of the Hilbert--Schmidt norm:
% ||A||_{L_2}^2 = sum_k ||A e_k||^2.

\begin{remark}
The key identity used above is $\sqrt{Q}\,e_k = \sqrt{\lambda_k}\,e_k$,
since $(e_k)$ are eigenvectors of $Q$ with eigenvalues $\lambda_k$,
and hence eigenvectors of $\sqrt{Q}$ with eigenvalues $\sqrt{\lambda_k}$.
\end{remark}

%% =============================
%% Summing over m to complete the isometry
%% =============================
\noindent
Summing over $m = 0, \ldots, k-1$:
\begin{align*}
    \sum_{m,n} E\!\bigl(\langle \Phi_m\,\Delta_m,\;
    \Phi_n\,\Delta_n\rangle_H\bigr)
    &\;=\;
    \sum_{m=0}^{k-1}
    E\!\Bigl(
        \bigl\|\Phi_m \circ \sqrt{Q}\,\bigr\|_{L_2(U,H)}^2
    \Bigr)
    \cdot (t_{m+1} - t_m) \\[6pt]
    &\;=\;
    E\!\left(
        \int_0^T
        \bigl\|\Phi_s \circ \sqrt{Q}\,\bigr\|_{L_2(U,H)}^2
        \,ds
    \right),
\end{align*}
which completes the proof of the Wiener--Itô isometry. \qed

\begin{remark}
The Wiener--It\^o isometry, once established for elementary processes, extends to all \textbf{predictable} integrands $\Phi$ satisfying $E\!\bigl(\int_0^T \|\Phi_s \circ \sqrt{Q}\|_{L_2(U,H)}^2\,ds\bigr) < \infty$.
The extension proceeds by a standard \textbf{completion argument}: approximate $\Phi$ by elementary processes $\Phi^{(n)}$ in the $L^2(\Omega \times [0,T])$ norm.  By the isometry, $\int \Phi^{(n)}\,dW$ is a Cauchy sequence in $\mathcal{M}^2$, and since $\mathcal{M}^2$ is complete (as a Banach space), the limit exists and defines $\int \Phi\,dW$.
This is completely analogous to how one extends the Lebesgue integral from simple functions to $L^1$, or the scalar It\^o integral from step processes to $L^2$-integrands.
\end{remark}

%% =============================
%% Adjoint operator and trace identity
%% =============================
\subsection*{Adjoint operator and trace identity}

% For A in L_2(U,H), there is a natural connection between
% the HS norm and the trace of A^*A.

Let $A \in L_2(U, H)$.  The \textbf{adjoint operator}
$A^* \colon H \to U$ satisfies $A^* \in L_2(H, U)$
(i.e.\ $A^*$ maps $H$ back to $U$).

\begin{remark}
One can prove that the composition $A^*\!A \colon U \to U$ is a
\textbf{finite-trace} operator on $U$.
% This is because A is Hilbert--Schmidt, so A^*A is trace class.
\end{remark}

\noindent
Indeed, its trace satisfies:
\begin{align*}
    \tr(A^*\!A)
    &\;=\;
    \sum_{k} \langle A^*\!A\,e_k,\; e_k \rangle_U \\[6pt]
    %% Use <A^*A e_k, e_k>_U = <A e_k, A e_k>_H by definition of adjoint.
    &\;=\;
    \sum_{k} \langle A\,e_k,\; A\,e_k \rangle_H \\[6pt]
    &\;=\;
    \sum_{k} \|A\,e_k\|_H^2 \\[6pt]
    &\;=\;
    \|A\|_{L_2(U,H)}^2
    \qquad \text{(HS norm)}.
\end{align*}

\begin{remark}
This identity shows that the Hilbert--Schmidt norm of $A$ equals
the square root of the trace of $A^*\!A$:
\[
    \|A\|_{L_2(U,H)}
    \;=\;
    \sqrt{\tr(A^*\!A)}.
\]
In particular, $A \in L_2(U,H)$ if and only if $A^*\!A$ is a
trace-class operator on $U$.
The key step uses the definition of the adjoint:
$\langle A^*\!A\,e_k, e_k \rangle_U = \langle A\,e_k, A\,e_k \rangle_H$,
which follows from $\langle A^*y, x \rangle_U = \langle y, Ax \rangle_H$
for all $x \in U$, $y \in H$.
\end{remark}

\end{document}
